\begingroup
\renewcommand{\arraystretch}{1.45}
\fancyhead[L]{\small LEMNISKATISCHE FUNKTIONEN II.}
\Needspace{14\baselineskip}
\CurvaArticle{4}{[4.]}
Werden die Werte $s\;[=\sin\frac{\pi}{\varpi}\varphi]$, die $\sin\operatorname{lemn}\varphi=0$ ergeben, nach den bekannten Regeln als Erzeuger eines unendlichen Produkts aufgefasst und ebenso die Werte von $s$, die $\sin\operatorname{lemn}\varphi=\infty$ ergeben, und sei das erste Produkt $P\varphi$, das zweite $Q\varphi$, so wird man (was wir streng beweisen können) setzen dürfen
\[
\sin\operatorname{lemn}\varphi=\dfrac{P\varphi}{Q\varphi}
\]


\Needspace{26\baselineskip}
\SourceLeaf{425}
\noindent Es ist dabei
\Needspace{12\baselineskip}
\[
\begin{aligned}P\varphi &=\alpha s
 \Bigl(1+\dfrac{4ss}{(e^{\pi}-e^{-\pi})^{2}}\Bigr)
 \Bigl(1+\dfrac{4ss}{(e^{2\pi}-e^{-2\pi})^{2}}\Bigr)
 \\[9pt]&\quad{}\Bigl(1+\dfrac{4ss}{(e^{3\pi}-e^{-3\pi})^{2}}\Bigr)\mathord{.}\mathord{.}\mathord{.}\end{aligned}
\]
\Needspace{12\baselineskip}
\[
\begin{aligned}Q\varphi &=
 \Bigl(1-\dfrac{4ss}{(e^{\frac12\pi}+e^{-\frac12\pi})^{2}}\Bigr)
 \Bigl(1-\dfrac{4ss}{(e^{\frac32\pi}+e^{-\frac32\pi})^{2}}\Bigr)
 \\[9pt]&\quad{}\Bigl(1-\dfrac{4ss}{(e^{\frac52\pi}+e^{-\frac52\pi})^{2}}\Bigr)\mathord{.}\mathord{.}\mathord{.}\end{aligned}
\]
\[
\begin{aligned}\alpha &=\dfrac{\varpi}{\pi}\end{aligned}
\]
\Needspace{28\baselineskip}
Setzt man ebenso
\Needspace{12\baselineskip}
\[
\begin{aligned}p\varphi &=c
 \Bigl(1-\dfrac{4ss}{(e^{\pi}+e^{-\pi})^{2}}\Bigr)
 \Bigl(1-\dfrac{4ss}{(e^{2\pi}+e^{-2\pi})^{2}}\Bigr)
 \\[9pt]&\quad{}\Bigl(1-\dfrac{4ss}{(e^{3\pi}+e^{-3\pi})^{2}}\Bigr)\mathord{.}\mathord{.}\mathord{.}\end{aligned}
\]
\Needspace{12\baselineskip}
\[
\begin{aligned}q\varphi &=
 \Bigl(1+\dfrac{4ss}{(e^{\frac12\pi}-e^{-\frac12\pi})^{2}}\Bigr)
 \Bigl(1+\dfrac{4ss}{(e^{\frac32\pi}-e^{-\frac32\pi})^{2}}\Bigr)
 \\[9pt]&\quad{}\Bigl(1+\dfrac{4ss'}{(e^{\frac52\pi}-e^{-\frac52\pi})^{2}}\Bigr)\mathord{.}\mathord{.}\mathord{.}\end{aligned}
\]
\noindent so ist
\[
\cos\operatorname{lemn}\varphi=\dfrac{p\varphi}{q\varphi}
\]
\Needspace{26\baselineskip}
\CurvaArticle{5}{[5.]}
\Needspace{12\baselineskip}
\[
\begin{aligned}
{}^{4}\!\surd 2\, .\,P(\varphi+\dfrac12\varpi)&=p\varphi\\[5pt]
{}^{4}\!\surd 2\, .\,Q(\varphi+\dfrac12\varpi)&=q\varphi
\end{aligned}
\]
\Needspace{12\baselineskip}
\[
\begin{aligned}
p(\varphi+\dfrac12\varpi)&=-{}^{4}\!\surd 2\, .\,P\varphi\\[5pt]
q(\varphi+\dfrac12\varpi)&=\phantom{-}{}^{4}\!\surd 2\, .\,Q\varphi
\end{aligned}
\]
\Needspace{16\baselineskip}
\[
\begin{aligned}
Pi\psi\varpi&=ie^{\pi\psi\psi}P\psi\varpi\\[5pt]
Qi\psi\varpi&=e^{\pi\psi\psi}Q\psi\varpi\\[5pt]
pi\psi\varpi&=e^{\pi\psi\psi}q\psi\varpi\\[5pt]
qi\psi\varpi&=e^{\pi\psi\psi}p\psi\varpi
\end{aligned}
\]
\Needspace{18\baselineskip}
\noindent [Werden die an anderer Stelle untersuchten Reihen auf Seite 405 dieses Bandes
\Needspace{16\baselineskip}
\[
\begin{aligned}
&\varphi-\dfrac1{60}\varphi^{5}-\dfrac1{10080}\varphi^{9}+\mathord{.}\mathord{.}\\[5pt]
&1+\dfrac1{12}\varphi^{4}-\dfrac1{10080}\varphi^{8}+\mathord{.}\mathord{.}\\[5pt]
&1-\dfrac12\varphi\varphi-\dfrac1{24}\varphi^{4}-\mathord{.}\mathord{.}\mathord{.}\\[5pt]
&1+\dfrac12\varphi\varphi-\dfrac1{24}\varphi^{4}+\mathord{.}\mathord{.}\mathord{.}
\end{aligned}
\]
\noindent der Reihe nach mit $\mathfrak{P}\varphi$, $\mathfrak{Q}\varphi$, $\mathfrak{p}\varphi$, $\mathfrak{q}\varphi$ bezeichnet, so gilt:]
\Needspace{16\baselineskip}
\[
\begin{aligned}
\mathfrak{P}\psi\varpi&=e^{\frac12\pi\psi\psi}P\psi\varpi
\\[5pt] \mathfrak{p}\psi\varpi&=e^{\frac12\pi\psi\psi}p\psi\varpi\\[5pt]
\mathfrak{Q}\psi\varpi&=e^{\frac12\pi\psi\psi}Q\psi\varpi
\\[5pt] \mathfrak{q}\psi\varpi&=e^{\frac12\pi\psi\psi}q\psi\varpi
\end{aligned}
\]



\Needspace{24\baselineskip}
\SourceLeaf{426}
\CurvaArticle{6}{[6.]}
Aus den oben angegebenen Ausdrücken folgt
\Needspace{22\baselineskip}
\[
\begin{aligned}
\sin\operatorname{lemn}\varphi
&=\dfrac{\dfrac{\pi}{\varpi}(e^{\frac12\pi}-e^{-\frac12\pi})}
        {e^{\frac12\pi}-2s+e^{-\frac12\pi}}
 -\dfrac{\dfrac{\pi}{\varpi}(e^{\frac12\pi}-e^{-\frac12\pi})}
        {e^{\frac12\pi}+2s+e^{-\frac12\pi}}\\[5pt]
&\quad-\dfrac{\dfrac{\pi}{\varpi}(e^{\frac32\pi}-e^{-\frac32\pi})}
        {e^{\frac32\pi}-2s+e^{-\frac32\pi}}
 +\dfrac{\dfrac{\pi}{\varpi}(e^{\frac32\pi}-e^{-\frac32\pi})}
        {e^{\frac32\pi}+2s+e^{-\frac32\pi}}\\[5pt]
&\quad+\text{usw.}\\[5pt]
\cos\operatorname{lemn}\varphi&= 
\end{aligned}
\]
\Needspace{14\baselineskip}
Hieraus folgt
\Needspace{14\baselineskip}
\[
\begin{aligned}\sin\operatorname{lemn}\psi\varpi
&=\dfrac{\pi}{\varpi}\dfrac{4}{e^{\frac12\pi}+e^{-\frac12\pi}}\sin\psi\pi
\\[7pt]&\quad{}-\dfrac{\pi}{\varpi}\dfrac{4}{e^{\frac32\pi}+e^{-\frac32\pi}}\sin3\psi\pi
+\mathord{.}\mathord{.}\mathord{.}\end{aligned}
\]
\clearpage
\endgroup
